% concatenated from maximaldilatation.tex
%%% ====================================================================
%%% @LaTeX-file{
%%%   filename  = "ijmsample.tex",
%%%   copyright = "Copyright 1995, 1999 American Mathematical Society,
%%%                2005 Hebrew University Magnes Press,
%%%                all rights reserved.  Copying of this file is
%%%                authorized only if either:
%%%                (1) you make absolutely no changes to your copy,
%%%                including name; OR
%%%                (2) if you do make changes, you first rename it
%%%                to some other name.",
%%% }
%%% ====================================================================
\NeedsTeXFormat{LaTeX2e}% LaTeX 2.09 can't be used (nor non-LaTeX)
[1994/12/01]% LaTeX date must December 1994 or later
%\documentclass{ijmart}
\documentclass{amsart}
%    Some definitions useful in producing this sort of documentation:
\chardef\bslash=`\\ % p. 424, TeXbook
%    Normalized (nonbold, nonitalic) tt font, to avoid font
%    substitution warning messages if tt is used inside section
%    headings and other places where odd font combinations might
%    result.

\usepackage{amsmath,amsfonts,amssymb}
\usepackage{mathabx}
\usepackage{graphicx}
\usepackage{url}
\usepackage{hyperref}
\usepackage{comment}
\usepackage{makecell}
\usepackage{boldline}
\setcellgapes{3pt}
\usepackage{hhline,colortbl}
\usepackage{xcolor}
\usepackage{textcomp}

\definecolor{blueviolet}{rgb}{0.5,0,1}


\newcommand{\ntt}{\normalfont\ttfamily}
%    command name
\newcommand{\cn}[1]{{\protect\ntt\bslash#1}}
%    LaTeX package name
\newcommand{\pkg}[1]{{\protect\ntt#1}}
%    File name
\newcommand{\fn}[1]{{\protect\ntt#1}}
%    environment name
\newcommand{\env}[1]{{\protect\ntt#1}}
\hfuzz1pc % Don't bother to report overfull boxes if overage is < 1pc

%       Theorem environments

%% \theoremstyle{plain} %% This is the default
\newtheorem{thm}{Theorem}[section]
\newtheorem{cor}[thm]{Corollary}
\newtheorem{lem}[thm]{Lemma}
\newtheorem{prop}[thm]{Proposition}
\newtheorem{con}[thm]{Conjecture}
\newtheorem{ex}[thm]{Example}
\newtheorem{ax}{Axiom}

\theoremstyle{definition}
\newtheorem{defn}{Definition}[section]
\newtheorem{rem}{Remark}[section]
\newtheorem*{notation}{Notation}

\theoremstyle{remark}
\newtheorem{step}{Step}

%\numberwithin{equation}{section}

\newcommand{\thmref}[1]{Theorem~\ref{#1}}
\newcommand{\secref}[1]{\S\ref{#1}}
\newcommand{\lemref}[1]{Lemma~\ref{#1}}


%       Math definitions

\newcommand{\A}{\mathcal{A}}
\newcommand{\B}{\mathcal{B}}
\newcommand{\st}{\sigma}
\newcommand{\XcY}{{(X,Y)}}
\newcommand{\SX}{{S_X}}
\newcommand{\SY}{{S_Y}}
\newcommand{\SXY}{{S_{X,Y}}}
\newcommand{\SXgYy}{{S_{X|Y}(y)}}
\newcommand{\Cw}[1]{{\hat C_#1(X|Y)}}
\newcommand{\G}{{G(X|Y)}}
\newcommand{\PY}{{P_{\mathcal{Y}}}}
\newcommand{\X}{\mathcal{X}}
\newcommand{\wt}{\widetilde}
\newcommand{\wh}{\widehat}
\newcommand{\Z}{\mathbb{Z}}

\DeclareMathOperator{\per}{per}
\DeclareMathOperator{\cov}{cov}
\DeclareMathOperator{\non}{non}
\DeclareMathOperator{\cf}{cf}
\DeclareMathOperator{\add}{add}
\DeclareMathOperator{\Cham}{Cham}
\DeclareMathOperator{\IM}{Im}
\DeclareMathOperator{\esssup}{ess\,sup}
\DeclareMathOperator{\meas}{meas}
\DeclareMathOperator{\seg}{seg}

\UseRawInputEncoding


%    \interval is used to provide better spacing after a [ that
%    is used as a closing delimiter.
\newcommand{\interval}[1]{\mathinner{#1}}

%    Notation for an expression evaluated at a particular condition. The
%    optional argument can be used to override automatic sizing of the
%    right vert bar, e.g. \eval[\biggr]{...}_{...}
\newcommand{\eval}[2][\right]{\relax
  \ifx#1\right\relax \left.\fi#2#1\rvert}

%    Enclose the argument in vert-bar delimiters:
\newcommand{\envert}[1]{\left\lvert#1\right\rvert}
\let\abs=\envert

%    Enclose the argument in double-vert-bar delimiters:
\newcommand{\enVert}[1]{\left\lVert#1\right\rVert}
\let\norm=\enVert

%\setcounter{tocdepth}{5}

\begin{document}
\title[Maximal dilatation]{Maximal dilatation on nonorientable surfaces\thanks{We thank Erwan Lanneau, Livio Liechti, Julien Marché, Alan Reid, and Darren Long.}}
\author[J. Ham]{Ji-Young Ham}
\address{Department of Mathematics, Kangwon National University, 1 Kangwondaehakgil, Chuncheon-si, Gangwon state, 24341\\
School of Liberal Arts, Seoul National University of Science and Technology, 232 Gongneung-ro, Nowon-gu, Seoul, 01811\\
  Korea} 
\email{jiyoungham@gmail.com}\thanks{The author is supported by Basic Science Research Program through the National Research Foundation of Korea (NRF) funded by the Ministry of Education, Science and Technology (No. NRF-2018005847).}

\author[J. Lee]{Joongul Lee}
\address{Department of Mathematics Education, Hongik University\\
94 Wausan-ro, Mapo-gu, Seoul 04066\\
   Korea}
\email{jglee@hongik.ac.kr\thanks{The author was supported by 2017 Hongik University Research Fund.}} 

\subjclass[2010]{37E30, 37B40, 57M60}

\keywords{Maximal degree, Nonorientable surface, Liechti-Strenner's polynomial, pseudo-Anosov, Minimal dilatation}

%\date{Received on MONTH, YEAR}
%\issueinfo{VOL}{NUM}{MONTH}{YEAR}
%\doiinfo{10.1007/DOI-NUMBER}
\begin{abstract}
 On each nonorientable surface of even genus $g \geq 4$, we show that the Liechti-Strenner's polynomial in~\cite{LiechtiStrenner18} gives a maximal dilatation among pseudo-Anosov diffeomorphisms with an orientable invariant foliation. This is proved by showing that this polynomial is irreducible.
 \end{abstract}
\maketitle
%\tableofcontents


\section{Introduction}
In 1970's, Thurston classified the mapping class group of a surface into periodic, pseudo-Anosov, and reducible~\cite{Thurston88}. The same content from a somewhat different point of view can be found in~\cite{CassonBleiler88}. In this paper, we adopt the Penner's approach~\cite{Penner88,Penner91}. Penner made use of bigon tracks, a slight generalization of train track. Nice example of bigon tracks can be found in~\cite{LiechtiStrenner18-1}. Our paper is based on 
Liechti-Strenner's pseudo-Anosov diffeomorphisms on nonorientable surfaces~\cite{LiechtiStrenner18}.

Let $\Sigma_g$ be a surface of finite type.
A diffeomorphism $h$ of $\Sigma_g$ is called \emph{pseudo-Anosov} if there is a pair of transversely measured foliations 
$\mathcal{F}^u$ and $\mathcal{F}^s$ in $\Sigma_g$ and a real number $\lambda >1$ such that $h(\mathcal{F}^u)=\lambda \mathcal{F}^u$ and  $h(\mathcal{F}^s)=1/\lambda \mathcal{F}^s$~\cite{Thurston88,CassonBleiler88}.
The number $\lambda$ is called the \emph{dilatation} of $h$ and the logarithm of $\lambda$ is called the \emph{topological entropy}. A beginner-friendly example about the dilatation can be found in~\cite{BaikRafiqiWu16}. The set of dilatations of pseudo-Anosov diffeomorphisms of the group of isotopy classes $\Sigma_g$ is discrete~\cite{ArnouxYoccoz81, Ivanov88}. In particular, there exists the minimal dilatation. Note that $\lambda$ is an algebraic number and hence there exists the minimal polynomial of it. Denote by the \emph{degree} of $\lambda$ the degree of the minimal polynomial of it. Denote by a \emph{maximal dilatation} on a set $S$ of diffeomorphisms of 
$\Sigma_g$ a dilatation that gives the maximal degree among $S$.

%%%%%%%%%%%%%%%%%%%%%%%%%%%%%%%%%%%%%%%%%%%%%%%%%%%%%%%%%%%

%%%%%%%%%%%%%%%%%%%%%%%%%%%%%%%%%%%%%%%%%%%%%%%%%%%%%%%%%%%%
 Liechti and Strenner~\cite{LiechtiStrenner18} determined the minimal dilatation of pseudo-Anosov diffeomorphisms with an orientable invariant foliation on the closed nonorientable surfaces of genus 4, 5, 6, 7, 8, 10, 12, 14, 16, 18 and 20. 
 Denote by $N_g$ the closed nonorientable surface of genus $g$ and by 
$\delta^{+} (N_g)$ the minimal dilatation among pseudo-Anosov diffeomorphisms of $N_g$ with an orientable invariant foliation. Liechti and Strenner conjectured that the largest root of $x^{2k-1}-x^k-x^{k-1}-1$ is the minimal dilatation of $N_{2k}$ for $k \geq 2$(Conjecture~\ref{con:1}). Denote by \emph{Liechti-Strenner’s polynomials} the odd degree polynomials in the conjecture~\ref{con:1}.
 
\begin{con}~\cite[Conjecture 1.2]{LiechtiStrenner18} \label{con:1}
For all $k \geq 2$, $\delta^{+} (N_{2k})$ is the largest root of $x^{2k-1}-x^k-x^{k-1}-1$.
\end{con}
%%%%%%%%%%%%%%%%%%%%%%%%%%%%%%%%%%%%%%%%%%%%%%%%%%%%%%%%%%%%

In~\cite{HamLee20}, we showed that the Liechti-Strenner's example for the closed nonorientable surface $N_{2k}$ in~\cite{LiechtiStrenner18} minimizes the dilatation within the class of pseudo-Anosov diffeomorphisms with an orientable invariant foliation and all but the first coefficient of the characteristic polynomial of the action induced on the first cohomology nonpositive. 



The main purpose of this paper is the following two theorems.

\begin{thm} \label{thm:irr}
Denote by $N_{2k}$ the closed nonorientable surface of genus $2k$. 
For all $k \geq 2$, the Liechti-Strenner\textquotesingle s polynomial,
$$ x^{2k-1}-x^k-x^{k-1}-1,$$ is irreducible.
\end{thm}

\begin{thm} \label{thm:main}
For all $k \geq 2$, the Liechti-Strenner\textquotesingle s polynomial,
$$ x^{2k-1}-x^k-x^{k-1}-1,$$ gives a maximal dilatation on $N_{2k}$ among pseudo-Anosov diffeomorphisms of $N_{2k}$ with an orientable invariant foliation.
\end{thm}




%%%%%%%%%%%%%%%%%%%%%%%%%%%%%%%%%%%%%%%%%%%%%%%%%%%%%%%%%%%%%
 It is known that this degree is possible~\cite{Strenner17}, but there have been no concrete examples of this being realized. The well known conjecture is that this polynomial gives the minimal dilatation~\cite{LiechtiStrenner18}. Here we prove the maximality instead of the minimailty. On each orientable surface of genus $g \geq 2$, a maximal dilatation among pseudo-Anosov diffeomorphisms with orientable invariant foliations is concretely presented in~\cite{LanneauLiechti24,Shin16}, but it doesn't seem to be the minimal dilatation. Earlier work concerning the degree of dilatations can be found in~\cite{ArnouxYoccoz81} and ~\cite{Long85}.
%%%%%%%%%%%%%%%%%%%%%%%%%%%%%%%%%%%%%%%%%%%%%%%%%%%%%%%%%%%%%%%%
Some references related to minimal dilatations are ~\cite{Abikoff80,AaberDunfield10,Bauer92,BestvinaHandel95,Birman74,Brinkmann00,Brinkmann04,ChoHam08,FarbMargalit12,FarberReinosoWang24,FLP79,Ham06,HamSong07,Hironaka10,HironakaKin06,KinTakasawa13,HironakaTsang24,Leininger04,LanneauThiffeault11,LanneauThiffeault11b,LiechtiStrenner18,LiechtiStrenner21,ErwanLivioCheuk25,Loving19,McMullen15,PapadopoulosPenner87,SnapPy,SongKoLos02,Tsai09,TsangZeng24,Valdivia12,Venzke08,Yazdi20,Zhirov95}.
The above references are neither exhaustive nor systematically assembled, so that their inclusion or noninclusion should not be taken as an indication of quality or relevance of the content. Nonetheless, we hope that you find them helpful.

%%%%%%%%%%%%%%%%%%%%%%%%%%%%%%%%%%%%%%%%%%%%%%%%%%%%%%%%%%%%
\section{Liechti-Strenner construction of nonorientable surfaces} \label{section:LS}

We will briefly introduce the Liechti-Strenner's method of constructing the nonorientable surface $\Sigma_{2k-1,2}$ of genus $g=2k$ and the Liechti-Strenner's polynomial.

\subsection{The graph $G_{2k-1,2}$}
Let $k \geq 2$ be an odd natural number. Let $G_{2k-1,2}$ be the graph whose vertices are the vertices of a regular $(2k-1)$-gon and every vertex $v$ is connected to the $2$ vertices that are the farthest away from $v$ in the cyclic order of the vertices. Figure~\ref{fig:graph} shows the graph 
$G_{3,2}$.

\begin{figure}
\begin{center}
\includegraphics[scale=1]{graph32-1.pdf}
\caption{The graph $G_{3,2}$.}
\label{fig:graph}
\end{center}
\end{figure}

\subsection{The surface $\Sigma_{2k-1,2}$}
For each $G_{2k-1,2}$, Liechti and Strenner constructed an nonorientable surface $\Sigma_{2k-1,2}$ that contains a collection of curves with
intersection graph $G_{2k-1,2}$~\cite[Subsection 2.2]{LiechtiStrenner18}. 
To construct $\Sigma_{2k-1,2}$, start with a disk with one crosscap. Next, we consider $4k-2$ disjoint intervals on the boundary of the disk and label
the intervals with integers from $1$ to $2k-1$ so that each label is used exactly twice.
In the cyclic order, the labels are $1,s,2,s+1,...,2k,s+2k$ where $s=k+2$ and
all labels are understood modulo $2k-1$.
For each label, the corresponding two intervals are connected by a twisted strip. Figure~\ref{fig:surface} shows the surface 
$\Sigma_{3,2}$.

\begin{figure}
\begin{center}
\includegraphics[scale=1]{sigma32-1}
\caption{The surface $\Sigma_{3,2}$ and the curve $c_1$.}
\label{fig:surface}
\end{center}
\end{figure}

\begin{lem}~\cite[Proposition 2.3]{LiechtiStrenner18}
The surface $\Sigma_{2k-1,2}$ is homeomorphic to the nonorientable surface of genus $2k$ with
$1$ boundary components.
\end{lem}

\subsection{The curves}
Liechti and Strenner constructed a two-sided curve $c_i$ for each label $i=1,\ldots,2k-1$ as follows. Each curve consists of two parts.
One part of each curve is the core of the strip corresponding to the label. The other part is an arc inside the disk that passes through the crosscap and connects the corresponding two intervals. The curve $c_1$ is shown in Figure~\ref{fig:surface}.
We choose markings for the $c_i$ which are invariant under the rotational symmetry See Figure~\ref{fig:curve}.

\begin{figure}
\begin{center}
\includegraphics[scale=1]{sigma32full-1}
\caption{A collection of filling inconsistently marked curves.}
\label{fig:curve}
\end{center}
\end{figure}

Note that every pair of curves intersects either once or not at all. The curves $c_i$ and $c_j$ are disjoint if and only if the two $i$ labels and the two $j$ labels link in the cyclic order. In other words, if the two $i$ labels separate the two $j$ labels.

\section {The mapping classes}

Denote by $r$ the rotation of $\Sigma_{2k-1,2}$ by one click in the clockwise direction. Denote by $T_{c_1}$ the right-handed Dehn twist about the curve $c_1$.
Define the mapping class
$$\Phi_k=r \circ T_{c_1}.$$
%$$\Phi_k=r^k  \circ T_{c_k}.$$
Note that $\Phi_k$ is pseudo-Anosov.


According to Penner~\cite[Theorem 3.1, Theorem 4.1]{Penner88} and Liechti and Strenner~\cite[Proposition 2.6]{LiechtiStrenner18}, we can construct a bigon track on 
$\Sigma_{2k-1,2}$ with filling curves $c_i$, $i=1,\ldots, 2k-1$. Each $c_i$ defines a characteristic measure $\mu_i$ on this bigon track, defined by assigning $1$ to the branches traversed by $c_i$ and zero to the rest. Let $H$ denote the cone generated by the measures $\mu_i$ in the cone of measures on the bigon track. 
Let $\mathcal{C}_k$ be $\mathcal{C}_k=\{c_i\,|i=1,\ldots,2k-1\}$ and $\mathcal{S}_k$ be the semigroup with presentation
$$\mathcal{S}_k (\mathcal{C}_k)=\langle c_i \in \mathcal{C}_k \, : \, c_i \leftrightarrow c_j \text{ if } c_i \cap c_j = \emptyset \rangle.$$

\vspace{2ex}

\begin{thm}~\cite[Theorem 3.4]{Penner88}
The action of $\mathcal{S}_k (\mathcal{C}_k)$ on $H$ admits a faithful representation as a semigroup of invertible (over $\Z$) positive matrices.
\end{thm}

Note that each $c_1$ corresponds to $T_{c_1}$ and the action of $c_1$ on $H$ in the basis $\mu_i$ is given by
$$I_{2k-1}+A$$
where $I_{2k-1}$ is the $(2k-1) \times (2k-1)$-identity matrix, and $A=\begin{bmatrix} a_{ij} \end{bmatrix}$ with $a_{ij}=0$ if $i \neq 1$, and $a_{1j}=\text{card}(c_1 \cap c_j)$.
The rotation $r$ acts by a permutation matrix.
The product of these two matrices is the companion matrix of the Liechti-Strenner's polynomial.

The following example shows how we obtain the Liechti-Strenner's polynomial on $\Sigma_{5,2}$.

\begin{ex}
On $\Sigma_{5,2}$, the action of $c_1$ on $H$ in the basis $\{\mu_i\}$ is given by

$$\begin{bmatrix}
1 & 0 & 1 & 1 & 0  \\
0 & 1 & 0 & 0 & 0  \\
0 & 0 & 1 & 0 & 0  \\
0 & 0 & 0 & 1 & 0  \\
0 & 0 & 0 & 0 & 1  \\

\end{bmatrix}.$$

The action of $r$ on $H$ in the basis $\mu_i$ is given by

$$\begin{bmatrix}
0 & 1 & 0 & 0 & 0 \\
0 & 0 & 1 & 0 & 0  \\
0 & 0 & 0 & 1 & 0 \\
0 & 0 & 0 & 0 & 1 \\
1 & 0 & 0 & 0 & 0  \\
\end{bmatrix}.$$

Hence $\Phi_3$ is represented by the below matrix $M_3$:

$$\begin{bmatrix}
0 & 1 & 0 & 0 & 0  \\
0 & 0 & 1 & 0 & 0  \\
0 & 0 & 0 & 1 & 0  \\
0 & 0 & 0 & 0 & 1  \\
1 & 0 & 1 & 1 & 0  \\
\end{bmatrix}.$$

The characteristic polynomial of $M_3$ is 
$$x^5-x^3-x^2-1.$$
\end{ex}

%%%%%%%%%%%%%%%%%%%%%%%%%%%%%%%%%%%%%%%%%%%%%%%%%%%%%%%%%%%%%%%
The proof of Theorem~\ref{thm:irr} is given in Section~\ref{sec:proof}. Since $x^3-x^2-x-1$ is irreducible mod 3, it is irreducible. The proof for $\Sigma_{5,2}$ and the proof for 
$\Sigma_{2k-1,2}$ are the essentially the same. We give the proof of $\Sigma_{5,2}$ as a warming up.

\begin{ex} \label{ex:2}
$x^5-x^3-x^2-1$ is irreducible.
\end{ex} 

\begin{proof}
$\bullet$ $x^5-x^3-x^2-1$ does not have linear factors.
\begin{enumerate} 
\item[I.] Suppose $x^5-x^3-x^2-1=(x^4+a_3 x^3+a_2 x^2+a_1 x+1)(b_1 x+b_0)$ where $b_0^2=1$ and $b_1^2=1$. 
Then, by comparing the coefficients, we have

%%%%%%%%%%%%%%%%%%%%
\begin{enumerate}
\item[(0)] $b_0=-1$ (the constants), \\
\item[(1)] $b_1+a_1 b_0=0$ (the coefficients of $x^1$), \\
\item[(2)] $a_1 b_1+a_2 b_0=-1$ (the coefficients of $x^2$) $\quad \Rightarrow \quad a_1 b_1+a_2 b_0+1=0$,\\
\item[(3)] $a_2 b_1+a_3 b_0=-1$ (the coefficients of $x^3$) $\quad \Rightarrow \quad a_2 b_1+a_3 b_0+1=0$, \\
\item[(4)] $a_3 b_1+b_0=0$ (the coefficients of $x^4$), \\
\item[(5)] $b_1=1$.
\end{enumerate}


We compute the resultant symmetrically in such a way that when we take the resultant of $(i)$ and $(j)$ with respect to $a_{t}$, we also take the resultant of $(5-i)$ and $(5-j)$ with respect to $a_{4-t}$.

\begin{enumerate}
\item[(6)] $res_{a_1}((1),(2))=-b_1^2+a_2b_0^2+b_0=-1+a_2+b_0$,
\item[(7)] $res_{a_3}((4),(3))=a_2 b_1^2+b_1-b_0^2=a_2+b_1-1$,
\end{enumerate}

Now,

\begin{equation*}
res_{a_2}((6),(7))=b_1-b_0
\end{equation*}

But then, since $b_1=1$, $b_0=1$ which is inconsistent with $(0)$.
%%%%%%%%%%%%%%%%%%%%%%%%%%%%%%%%%%%%%%%%%%%%%%%%%%
\item[II.] Suppose $x^5-x^3-x^2-1=(a_4x^4+a_3 x^3+a_2 x^2+a_1 x+a_0)(x+1)$ where $a_0^2=1$ and $a_4^2=1$. 
Then, by comparing the coefficients, we have

%%%%%%%%%%%%%%%%%%%%
\begin{enumerate}
\item[(0)] $a_0=-1$ (the constants), \\
\item[(1)] $a_1+ a_0=0$ (the coefficients of $x^1$), \\
\item[(2)] $a_1+a_2=-1$ (the coefficients of $x^2$) $\quad \Rightarrow \quad a_1+a_2+1=0$,\\
\item[(3)] $a_2+a_3=-1$ (the coefficients of $x^3$) $\quad \Rightarrow \quad a_2+a_3+1=0$, \\
\item[(4)] $a_3+a_4=0$ (the coefficients of $x^4$), \\
\item[(5)] $a_4=1$.
\end{enumerate}


We compute the resultant symmetrically in such a way that when we take the resultant of $(i)$ and $(j)$ with respect to $a_{t}$, we also take the resultant of $(5-i)$ and $(5-j)$ with respect to $a_{4-t}$.

\begin{enumerate}
\item[(6)] $res_{a_1}((1),(2))=a_2-a_0+1$,
\item[(7)] $res_{a_3}((4),(3))=a_2-a_4+1$.
\end{enumerate}

Now,

\begin{equation*}
res_{a_2}((6),(7))=a_0-a_4
\end{equation*}

But then, since $a_4=1$, $a_0=1$ which is inconsistent with $(0)$.

\end{enumerate} 
%%%%%%%%%%%%%%%%%%%%%%%%%%%%%%%%%%%%%%%%%%%%%%%%%%%%%%

$\bullet$ $x^5-x^3-x^2-1$ does not have quadratic factors.
\begin{enumerate} 
\item[I.] Suppose $x^5-x^3-x^2-1=(x^3+a_2 x^2+a_1 x+1)(b_2 x^2+b_1 x+b_0)$ where $b_0^2=1$ and $b_2^2=1$.
Similarly, we symmetrically take the resultant several times until we get two equations with three variables $b_2$,$b_1$, and $b_0$.
Now, taking the resultant one more time with respect to $b_1$, we get
 \begin{equation*}
 4(1-b_0 b_2)=4b_2(b_2-b_0) .
\end{equation*}
But then, the system becomes inconsistent.
\item[II.] Suppose $x^5-x^3-x^2-1=(a_3 x^3+a_2 x^2+a_1 x+a_0)(x^2+b_1 x+1)$.
Similarly, we have $2(1-a_0 a_3)=2a_3(a_3-a_0).$
But then, the system becomes inconsistent.
\end{enumerate}

Therefore, by two bullets, $x^5-x^3-x^2-1$ is irreducible.
\end{proof}
%%%%%%%%%%%%%%%%%%%%%%%%%%%%%%%%%%%%%%%%%%%%%%%%%%%%%%

%%%%%%%%%%%%%%%%%%%%%%%%%%%%%%%%%%%%%%%%%%%%%%%%%%%%%%%%%%%%%%%%%%%%%%%%%%%%%%%%%%%%%%%%%%%%%%%%%%%%%%%%%%%%%%%%%%%%%%%%%
\section{proof of Theorem~\ref{thm:irr}} \label{sec:proof}

We may assume that $k > 3$.

Suppose $$x^{2k-1}-x^k-x^{k-1}-1=  \left(\sum_{i=0}^{l} a_i x^i \right)  \left(\sum_{j=0}^{m} b_j x^j \right)$$ where $l+m=2k-1$.
Without loss of generality, we assume that $l$ is even. We assume that $b_0=b_m=1$  and $a_0^2=a_l^2=1$ (or $a_0=a_l=1$ and $b_0^2=b_l^2=1$). Then, we symmetrically take the resultant several times until we get two equations with three variables, $a_l$, $a_{l/2}$, and $a_0$ ( or $b_m$, $a_{l/2}$, and $b_0$). 
Now taking the resultant one more time with respect to $a_{l/2}$, we have an equation of $a_l$ and $a_0$ (or $b_m$ and $b_0$). We took the resultant so that the coefficients of $a_l$ and $a_0$ (or $b_m$ and $b_0$) are the same η in the final polynomial,  but have different signs. 
If the $\eta \neq 0$, we get an inconsistent system.  At each step, we carefully choose a pair of polynomials so that the two polynomials in the final resultant are not equal up to the powers of $a_l$ and $a_0$ (or $b_m$ and $b_0$). This ensures $\eta \neq 0$. 
Reducing variables through substitution before taking the resultants, greatly reduces work.  Also, the lower the degree, the fewer fakes there are.
%%%%%%%%%%%%%%%%%%%%%%%%%%%%%%%%%%%%%%%%%%%%%%%%%%%%%%%%%%%%%%%

\section{proof of the Theorem~\ref{thm:main} on $\Sigma_{2k-1,2}$}

Since $\text{dim}(H_1(N_{2k},\mathbb{R}))=2k-1$, Theorem~\ref{thm:main} is proved by Theorem~\ref{thm:irr}.
\section{acknowledgement}
We thank Erwan Lanneau, Livio Liechti, Julien Marché, Alan Reid, and Darren Long. 
This work was supported by Basic Science Research Program through the National Research Foundation of Korea (NRF) funded by the Ministry of Education, Science and Technology (No. NRF-2018005847). %The second author was supported by 2017 Hongik University Research Fund.
%%%%%%%%%%%%%%%%%%%%%%%%%%%%%%%%%%%%%%%%%%%%%%%%%%%%%%%%%%%%%%%%%%%
%%%%%%%%%%%%%%%%%%%%%%%%%%%%%%%%%%%%%%%%%%%%%%%%%%%%%%%%%%%%%%%%%%%%%%
%%%%%%%%%%%%%%%%%%%%%%%%%%%%%%%%%%%%%%%%%%%%%%%%%%%%%%%%%%%%
%\bibliographystyle{amsalpha}
%\bibliography{minimalreference}
\providecommand{\bysame}{\leavevmode\hbox to3em{\hrulefill}\thinspace}
\providecommand{\MR}{\relax\ifhmode\unskip\space\fi MR }
% \MRhref is called by the amsart/book/proc definition of \MR.
\providecommand{\MRhref}[2]{%
  \href{http://www.ams.org/mathscinet-getitem?mr=#1}{#2}
}
\providecommand{\href}[2]{#2}
\begin{thebibliography}{BRW16}

\bibitem[Abi80]{Abikoff80}
William Abikoff, \emph{The real analytic theory of {T}eichm\"uller space},
  Lecture Notes in Mathematics, vol. 820, Springer, Berlin, 1980. \MR{590044}

\bibitem[AD10]{AaberDunfield10}
John~W. Aaber and Nathan Dunfield, \emph{Closed surface bundles of least
  volume}, Algebr. Geom. Topol. \textbf{10} (2010), no.~4, 2315--2342.
  \MR{2745673}

\bibitem[AY81]{ArnouxYoccoz81}
Pierre Arnoux and Jean-Christophe Yoccoz, \emph{Construction de
  diff\'{e}omorphismes pseudo-{A}nosov}, C. R. Acad. Sci. Paris S\'{e}r. I
  Math. \textbf{292} (1981), no.~1, 75--78. \MR{610152}

\bibitem[Bau92]{Bauer92}
Max Bauer, \emph{An upper bound for the least dilatation}, Trans. Amer. Math.
  Soc. \textbf{330} (1992), no.~1, 361--370. \MR{1094556}

\bibitem[BH95]{BestvinaHandel95}
M.~Bestvina and M.~Handel, \emph{Train-tracks for surface homeomorphisms},
  Topology \textbf{34} (1995), no.~1, 109--140. \MR{1308491}

\bibitem[Bir74]{Birman74}
Joan~S. Birman, \emph{Braids, links, and mapping class groups}, Princeton
  University Press, Princeton, N.J.; University of Tokyo Press, Tokyo, 1974,
  Annals of Mathematics Studies, No. 82. \MR{0375281}

\bibitem[Bri00]{Brinkmann00}
Peter Brinkmann, \emph{An implementation of the {B}estvina-{H}andel algorithm
  for surface homeomorphisms}, Experiment. Math. \textbf{9} (2000), no.~2,
  235--240. \MR{1780208}

\bibitem[Bri04]{Brinkmann04}
\bysame, \emph{A note on pseudo-{A}nosov maps with small growth rate},
  Experiment. Math. \textbf{13} (2004), no.~1, 49--53. \MR{2065567}

\bibitem[BRW16]{BaikRafiqiWu16}
Hyungryul Baik, Ahmad Rafiqi, and Chenxi Wu, \emph{Constructing pseudo-{A}nosov
  maps with given dilatations}, Geom. Dedicata \textbf{180} (2016), 39--48.
  \MR{3451455}

\bibitem[CB88]{CassonBleiler88}
Andrew~J. Casson and Steven~A. Bleiler, \emph{Automorphisms of surfaces after
  {N}ielsen and {T}hurston}, London Mathematical Society Student Texts, vol.~9,
  Cambridge University Press, Cambridge, 1988. \MR{964685}

\bibitem[CDWo]{SnapPy}
Marc Culler, Nathan Dunfield, Jeff Weeks, and Many others, \emph{{S}nap{P}y},
  \url{http://www.math.uic.edu/t3m/SnapPy/}.

\bibitem[CH08]{ChoHam08}
Jin-Hwan Cho and Ji-Young Ham, \emph{The minimal dilatation of a genus-two
  surface}, Experiment. Math. \textbf{17} (2008), no.~3, 257--267. \MR{2455699}

\bibitem[FLPÌ79]{FLP79}
A.~Fathi, F.~Laudenbach, and V.~Po~́enaru, \emph{Travaux de {T}hurston sur les
  surfaces}, Ast\'erisque, vol. 66-67, Soci\'et\'e{} Math\'ematique de France,
  Paris, 1979, S\'eminaire Orsay, With an English summary. \MR{568308}

\bibitem[FM12]{FarbMargalit12}
Benson Farb and Dan Margalit, \emph{A primer on mapping class groups},
  Princeton Mathematical Series, vol.~49, Princeton University Press,
  Princeton, NJ, 2012. \MR{2850125}

\bibitem[FRW24]{FarberReinosoWang24}
Ethan Farber, Braeden Reinoso, and Luya Wang, \emph{Fixed-point-free
  pseudo-{A}nosov homeomorphisms, knot {F}loer homology and the cinquefoil},
  Geom. Topol. \textbf{28} (2024), no.~9, 4337--4381. \MR{4845456}

\bibitem[Ham06]{Ham06}
Ji-young Ham, \emph{The minimal dilatations of 4 and 5 braids}, ProQuest LLC,
  Ann Arbor, MI, 2006, Thesis (Ph.D.)--University of California, Santa Barbara.
  \MR{2708350}

\bibitem[Hir10]{Hironaka10}
Eriko Hironaka, \emph{Small dilatation mapping classes coming from the simplest
  hyperbolic braid}, Algebr. Geom. Topol. \textbf{10} (2010), no.~4,
  2041--2060. \MR{2728483}

\bibitem[HK06]{HironakaKin06}
Eriko Hironaka and Eiko Kin, \emph{A family of pseudo-{A}nosov braids with
  small dilatation}, Algebr. Geom. Topol. \textbf{6} (2006), 699--738.
  \MR{2240913}

\bibitem[HL20]{HamLee20}
Ji-Young Ham and Joongul Lee, \emph{Remarks on the {L}iechti-{S}trenner's
  examples having small dilatations}, Commun. Korean Math. Soc. \textbf{35}
  (2020), no.~4, 1299--1307. \MR{4169009}

\bibitem[HS07]{HamSong07}
Ji-Young Ham and Won~Taek Song, \emph{The minimum dilatation of pseudo-{A}nosov
  5-braids}, Experiment. Math. \textbf{16} (2007), no.~2, 167--179.
  \MR{2339273}

\bibitem[HT24]{HironakaTsang24}
Eriko Hironaka and Chi~Cheuk Tsang, \emph{Standardly embedded train tracks and
  pseudo-anosov maps with minimum expansion factor}, \url{arXiv:2210.13418},
  2024, Groups, Geometry, and Dynamics (to appear).

\bibitem[Iva88]{Ivanov88}
N.~V. Ivanov, \emph{Coefficients of expansion of pseudo-{A}nosov
  homeomorphisms}, Zap. Nauchn. Sem. Leningrad. Otdel. Mat. Inst. Steklov.
  (LOMI) \textbf{167} (1988), no.~Issled. Topol. 6, 111--116, 191. \MR{964259}

\bibitem[KT13]{KinTakasawa13}
Eiko Kin and Mitsuhiko Takasawa, \emph{Pseudo-{A}nosovs on closed surfaces
  having small entropy and the {W}hitehead sister link exterior}, J. Math. Soc.
  Japan \textbf{65} (2013), no.~2, 411--446. \MR{3055592}

\bibitem[Lei04]{Leininger04}
Christopher~J. Leininger, \emph{On groups generated by two positive
  multi-twists: {T}eichm\"{u}ller curves and {L}ehmer's number}, Geom. Topol.
  \textbf{8} (2004), 1301--1359. \MR{2119298}

\bibitem[LL24]{LanneauLiechti24}
Erwan Lanneau and Livio Liechti, \emph{Trace field degrees of {A}belian
  differentials}, Comment. Math. Helv. \textbf{99} (2024), no.~1, 81--110.
  \MR{4709307}

\bibitem[LLT25]{ErwanLivioCheuk25}
Erwan Lanneau, Livio Liechti, and Chi~Cheuk Tsang, \emph{Minimal stretch
  factors of orientation-reversing fully-punctured pseudo-{A}nosov maps},
  Trans. Amer. Math. Soc. \textbf{378} (2025), no.~4, 2943--2968. \MR{4880466}

\bibitem[Lon85]{Long85}
D.~D. Long, \emph{Constructing pseudo-{A}nosov maps}, Knot theory and manifolds
  ({V}ancouver, {B}.{C}., 1983), Lecture Notes in Math., vol. 1144, Springer,
  Berlin, 1985, pp.~108--114. \MR{823284}

\bibitem[Lov19]{Loving19}
Marissa Loving, \emph{Least dilatation of pure surface braids}, Algebr. Geom.
  Topol. \textbf{19} (2019), no.~2, 941--964, Appendix written jointly with
  Hugo Parlier. \MR{3924180}

\bibitem[LS20a]{LiechtiStrenner18-1}
Livio Liechti and Bal\'azs Strenner, \emph{The {A}rnoux-{Y}occoz mapping
  classes via {P}enner's construction}, Bull. Soc. Math. France \textbf{148}
  (2020), no.~3, 383--397. \MR{4161163}

\bibitem[LS20b]{LiechtiStrenner18}
\bysame, \emph{Minimal pseudo-{A}nosov stretch factors on nonoriented
  surfaces}, Algebr. Geom. Topol. \textbf{20} (2020), no.~1, 451--485.
  \MR{4071380}

\bibitem[LS21]{LiechtiStrenner21}
\bysame, \emph{Minimal {P}enner dilatations on nonorientable surfaces}, J.
  Topol. Anal. \textbf{13} (2021), no.~1, 187--218. \MR{4243077}

\bibitem[LT11a]{LanneauThiffeault11}
Erwan Lanneau and Jean-Luc Thiffeault, \emph{On the minimum dilatation of
  braids on punctured discs}, Geom. Dedicata \textbf{152} (2011), 165--182,
  Supplementary material available online. \MR{2795241}

\bibitem[LT11b]{LanneauThiffeault11b}
\bysame, \emph{On the minimum dilatation of pseudo-{A}nosov homeromorphisms on
  surfaces of small genus}, Ann. Inst. Fourier (Grenoble) \textbf{61} (2011),
  no.~1, 105--144. \MR{2828128}

\bibitem[McM15]{McMullen15}
Curtis~T. McMullen, \emph{Entropy and the clique polynomial}, J. Topol.
  \textbf{8} (2015), no.~1, 184--212. \MR{3335252}

\bibitem[Pen88]{Penner88}
Robert~C. Penner, \emph{A construction of pseudo-{A}nosov homeomorphisms},
  Trans. Amer. Math. Soc. \textbf{310} (1988), no.~1, 179--197. \MR{930079}

\bibitem[Pen91]{Penner91}
R.~C. Penner, \emph{Bounds on least dilatations}, Proc. Amer. Math. Soc.
  \textbf{113} (1991), no.~2, 443--450. \MR{1068128}

\bibitem[PP87]{PapadopoulosPenner87}
Athanase Papadopoulos and Robert~C. Penner, \emph{A characterization of
  pseudo-{A}nosov foliations}, Pacific J. Math. \textbf{130} (1987), no.~2,
  359--377. \MR{914107}

\bibitem[Shi16]{Shin16}
Hyunshik Shin, \emph{Algebraic degrees of stretch factors in mapping class
  groups}, Algebr. Geom. Topol. \textbf{16} (2016), no.~3, 1567--1584.
  \MR{3523049}

\bibitem[SKL02]{SongKoLos02}
Won~Taek Song, Ki~Hyoung Ko, and J\'er\^ome~E. Los, \emph{Entropies of braids},
  vol.~11, 2002, Knots 2000 Korea, Vol. 2 (Yongpyong), pp.~647--666.
  \MR{1915500}

\bibitem[Str17]{Strenner17}
Bal\'azs Strenner, \emph{Algebraic degrees of pseudo-{A}nosov stretch factors},
  Geom. Funct. Anal. \textbf{27} (2017), no.~6, 1497--1539. \MR{3737368}

\bibitem[Thu88]{Thurston88}
William~P. Thurston, \emph{On the geometry and dynamics of diffeomorphisms of
  surfaces}, Bull. Amer. Math. Soc. (N.S.) \textbf{19} (1988), no.~2, 417--431.
  \MR{956596}

\bibitem[Tsa09]{Tsai09}
Chia-Yen Tsai, \emph{The asymptotic behavior of least pseudo-{A}nosov
  dilatations}, Geom. Topol. \textbf{13} (2009), no.~4, 2253--2278.
  \MR{2507119}

\bibitem[TZ24]{TsangZeng24}
CHI~CHEUK TSANG and XIANGZHUO ZENG, \emph{Minimum dilatations of pseudo-anosov
  braids}, \url{arXiv:2412.01648}, 2024, Preprint.

\bibitem[Val12]{Valdivia12}
Aaron~D. Valdivia, \emph{Sequences of pseudo-{A}nosov mapping classes and their
  asymptotic behavior}, New York J. Math. \textbf{18} (2012), 609--620.
  \MR{2967106}

\bibitem[Ven08]{Venzke08}
Rupert~William Venzke, \emph{Braid {F}orcing, {H}yperbolic {G}eometry, and
  {P}seudo-{A}nosov {S}equences of {L}ow {E}ntropy}, ProQuest LLC, Ann Arbor,
  MI, 2008, Thesis (Ph.D.)--California Institute of Technology. \MR{3078552}

\bibitem[Yaz20]{Yazdi20}
Mehdi Yazdi, \emph{Pseudo-{A}nosov maps with small stretch factors on punctured
  surfaces}, Algebr. Geom. Topol. \textbf{20} (2020), no.~4, 2095--2128.
  \MR{4127091}

\bibitem[Zhi95]{Zhirov95}
A.~Yu. Zhirov, \emph{On the minimum dilation of pseudo-{A}nosov diffeomorphisms
  of a double torus}, Uspekhi Mat. Nauk \textbf{50} (1995), no.~1(301),
  197--198. \MR{1331364}

\end{thebibliography}

%%%%%%%%%%%%%%%%%%%%%%%%%%%%%%%%%%%%%%%%%%%%%%%%%%%%%%%%%%%%
\end{document}
\endinput
